% ============================================================
%  Expectation Maximisers — Report Sketch
%  Fill in every \todo{} block before submission.
% ============================================================
\documentclass[11pt, a4paper]{article}

% ── Packages ────────────────────────────────────────────────
\usepackage[a4paper, margin=1in]{geometry}
\usepackage{amsmath, amssymb}
\usepackage{booktabs}
\usepackage{array}
\usepackage{hyperref}
\usepackage{xcolor}
\usepackage{graphicx}
\usepackage{listings}
\usepackage{enumitem}
\usepackage{titlesec}
\usepackage{parskip}        % space between paragraphs, no indent
% microtype omitted for compatibility
\usepackage{caption}
\usepackage{float}

% ── Colours ─────────────────────────────────────────────────
\definecolor{todocolor}{RGB}{192, 57, 43}
\definecolor{warncolor}{RGB}{180, 95, 6}
\definecolor{codeblue}{RGB}{31, 56, 100}
\definecolor{codegray}{RGB}{245, 245, 245}

% ── TODO / CAUTION commands ──────────────────────────────────
\newcommand{\todo}[1]{%
  \textcolor{todocolor}{\textbf{[TODO:} \textit{#1}\textbf{]}}%
}
% \warn{} flags a design or correctness concern that needs resolving before the demo.
\newcommand{\warn}[1]{%
  \noindent\colorbox{yellow!25}{\parbox{\dimexpr\linewidth-2\fboxsep}{%
    \textcolor{warncolor}{\textbf{$\triangleright$ CAUTION:} \textit{#1}}%
  }}%
  \vspace{4pt}%
}

% ── Hyperref setup ──────────────────────────────────────────
\hypersetup{
  colorlinks=true,
  linkcolor=codeblue,
  urlcolor=codeblue,
  citecolor=codeblue,
  pdftitle={Expectation Maximisers},
  pdfauthor={},
}

% ── Code listing style ──────────────────────────────────────
\lstset{
  backgroundcolor=\color{codegray},
  basicstyle=\ttfamily\small,
  breaklines=true,
  frame=single,
  rulecolor=\color{gray!40},
  tabsize=2,
  showstringspaces=false,
}

% ── Section title formatting ────────────────────────────────
\titleformat{\section}{\large\bfseries\color{codeblue}}{{\thesection}}{1em}{}
\titleformat{\subsection}{\normalsize\bfseries\color{codeblue}}{\thesubsection}{1em}{}
\titleformat{\subsubsection}{\normalsize\bfseries}{\thesubsubsection}{1em}{}

% ============================================================
\begin{document}

% ── Title page ──────────────────────────────────────────────
\begin{titlepage}
  \centering
  \vspace*{3cm}
  {\Huge\bfseries\color{codeblue} Expectation Maximisers \par}
  \vspace{0.6cm}
  {\large\itshape Crowd Preference Learning with EM-Weighted Preference Optimisation \par}
  \vspace{1.5cm}
  \rule{\linewidth}{0.4pt}\\[0.4cm]
  % {\large Course Project Report \par}
  \vspace{0.4cm}
  \todo{Author names} \par
  \vspace{0.2cm}
  CS 6103 - Human Centred AI\par
  \vspace{0.2cm}
  {\large May 2026 \par}
  \vspace{1cm}
  \rule{\linewidth}{0.4pt}
\end{titlepage}

% ── Table of contents ───────────────────────────────────────
\tableofcontents
\newpage

% ============================================================
\section{Abstract}

This project combines noisy crowd annotations on the Anthropic HH-RLHF dataset with
Dawid--Skene-style Expectation Maximization (EM) and EM-weighted DPO-style preference
optimisation. Annotators vote on which of two responses (A or B) is better; votes are
simulated with heterogeneous worker reliability. EM infers per-prompt latent preferences
and per-worker accuracy parameters. A policy language model (GPT-2) supplies LLM priors
for the EM E-step; training uses a weighted objective where soft labels come from EM's
posterior. The GPT-2 pipeline achieves approximately 75\% golden-dataset accuracy.

\todo{Add 2--3 sentences on key findings / contributions once experiments are complete.}

\newpage

% ============================================================
\section{Introduction}

\subsection{Motivation}

Human feedback is central to aligning language models with human preferences. In practice,
annotator disagreement is the norm: workers vary widely in expertise, attention, and
subjective judgement. Treating all votes equally discards valuable signal about annotator
reliability.

\todo{Expand with literature context: RLHF, Bradley--Terry models, disagreement in annotation.}

\subsection{Problem Statement}

Given a set of pairwise preference votes from heterogeneous annotators on prompt--response
pairs, we aim to (1) infer the latent ground-truth preference per prompt and (2) train a
language model policy using these inferred preferences as soft labels.

\subsection{Contributions}

\begin{itemize}[noitemsep]
  \item A simulation framework over Anthropic HH-RLHF with four annotator types
        (expert, average, spammer, adversary) and a bipartite assignment graph.
  \item A staged experimental pipeline: plain DPO baseline $\to$ offline EM + DPO
        $\to$ offline EM + projected DPO + SFT, isolating the contribution of each component.
  \item An EM quality metric ($\mathcal{L}_\text{EM}$) combining label MSE and worker
        parameter MSE, used to track EM convergence visually across phases.
  \item Stress-testing EM under increasing task scale, adversarial annotator ratios, and
        vote sparsity, with individual and combined factor analyses.
  \item A robustness evaluation showing that EM + best DPO outperforms best DPO alone
        even on the worst-case annotator regime.
  \item (If time) Hierarchical EM with annotator expertise levels, followed by projected
        DPO + SFT for peak accuracy.
\end{itemize}

\todo{Revise after finalising experiments.}

\newpage

% ============================================================
\section{Background \& Related Work}

\subsection{Reinforcement Learning from Human Feedback (RLHF)}

\todo{Summarise the RLHF pipeline: SFT $\to$ reward model $\to$ PPO.
Cite Christiano et al.~(2017), Stiennon et al.~(2020), InstructGPT.}

\subsection{Direct Preference Optimisation (DPO)}

\todo{Describe the DPO objective (Rafailov et al.~2023): avoids an explicit reward model
by reparameterising as policy log-ratios.}

\subsection{The Dawid--Skene Model}

The Dawid--Skene model \cite{dawid1979} is an EM algorithm for inferring latent
ground-truth labels from noisy crowd annotations. Each annotator is characterised by a
confusion matrix; in the binary case this reduces to per-worker sensitivity $\alpha_j$
and specificity $\beta_j$.

\todo{Add 1--2 sentences on recent extensions and applications to NLP annotation.}

\subsection{Dataset: Anthropic HH-RLHF}

Anthropic's Helpful and Harmless RLHF dataset \cite{bai2022} contains human-written
preference pairs (chosen / rejected responses). This project uses the first 1,000 rows
of the training split as the prompt pool.

\newpage

% ============================================================
\section{Method}

\subsection{Simulation}

A bipartite graph assigns 1,000 prompts to 50 annotators (3 workers per prompt,
60 prompts per worker). Table~\ref{tab:workers} describes the four worker types.

\begin{table}[H]
  \centering
  \caption{Simulated annotator types and their properties.}
  \label{tab:workers}
  \begin{tabular}{llll}
    \toprule
    \textbf{Type} & \textbf{Count} & \textbf{True Accuracy} & \textbf{Behaviour} \\
    \midrule
    Expert    & 10 & High    & Reliably votes with ground truth \\
    Average   & 25 & Medium  & Moderately reliable \\
    Spammer   & 10 & $\approx$0.5 & Near-random votes \\
    Adversary &  5 & Low     & Systematically votes against truth \\
    \bottomrule
  \end{tabular}
\end{table}

For each prompt, \texttt{truth\_is\_A} is drawn uniformly at random. Response A is the
HH-RLHF \textit{chosen} response when \texttt{truth\_is\_A}$=1$, else \textit{rejected}.
Each vote is sampled as:
\[
  v_{ij} = \begin{cases} \text{truth}_i & \text{with prob. } \theta_j \\ 1-\text{truth}_i & \text{otherwise} \end{cases}
\]

\subsection{Expectation Maximisation (Dawid--Skene)}

\subsubsection{Latent Variables}

Let $\gamma_i = P(\text{prompt } i \text{'s true label is ``A is better''})$, with
per-annotator parameters $\alpha_j = P(\text{vote A} \mid \text{truth A})$ and
$\beta_j = P(\text{vote B} \mid \text{truth B})$. The global prior is
$\pi_1 = \frac{1}{N}\sum_i \gamma_i$.

\subsubsection{E-Step}

The posterior over the latent label for prompt $i$ is accumulated in log-space over
all votes from assigned workers:

\[
  \log \gamma_i \propto \log \pi_1
    + \sum_{j \in \mathcal{A}(i)} \bigl[v_{ij}\log\alpha_j + (1-v_{ij})\log(1-\beta_j)\bigr]
\]

When LLM priors $p_i^\text{LLM}$ are available (see \S\ref{sec:joint}), the term
$\log\pi_1$ is replaced by $\log p_i^\text{LLM}$.

\todo{Add the symmetric expression for $\log(1-\gamma_i)$ and normalisation step.}

\subsubsection{M-Step}

Given the updated $\gamma_i$, expected counts are accumulated via \texttt{scatter\_add}
and used to update $\pi_1$, $\alpha_j$, $\beta_j$:

\[
  \alpha_j = \frac{\sum_{i: j\in\mathcal{A}(i)} \gamma_i \cdot \mathbf{1}[v_{ij}=1]}
                  {\sum_{i: j\in\mathcal{A}(i)} \gamma_i}, \qquad
  \beta_j  = \frac{\sum_{i: j\in\mathcal{A}(i)} (1-\gamma_i) \cdot \mathbf{1}[v_{ij}=0]}
                  {\sum_{i: j\in\mathcal{A}(i)} (1-\gamma_i)}
\]

$\alpha_j$ and $\beta_j$ are clamped away from 0 and 1 for numerical stability.

\subsection{Weighted DPO-Style Loss}

The implicit reward for a sequence $y$ under policy $\pi_\theta$ and reference $\pi_\text{ref}$ is:
\[
  r(x, y) = \beta \left(\log\frac{\pi_\theta(y\mid x)}{\pi_\text{ref}(y\mid x)}\right)
\]

The weighted loss for prompt $i$ with soft label $\gamma_i$ is:
\[
  \mathcal{L}_i = -\,\gamma_i \log\sigma\!\bigl(r(x_i, y_i^A) - r(x_i, y_i^B)\bigr)
                 -(1-\gamma_i)\log\sigma\!\bigl(r(x_i, y_i^B) - r(x_i, y_i^A)\bigr)
\]

where $y_i^A$ and $y_i^B$ are the two candidate responses.

\subsection{Joint Training Loop}
\label{sec:joint}

Each training epoch alternates three phases:

\begin{enumerate}[noitemsep]
  \item \textbf{LLM Priors.} Compute per-prompt
        $p_i^\text{LLM} = \sigma\!\bigl(\beta(r(x_i,y_i^A)-r(x_i,y_i^B))\bigr)$
        under the current policy and frozen reference.
  \item \textbf{EM Update.} Run one E-step (with $p_i^\text{LLM}$ as prior) and one
        M-step to update $\gamma_i$, $\alpha_j$, $\beta_j$.
  \item \textbf{Policy Update.} Compute $\mathcal{L}_i$ using the updated $\gamma_i$;
        backpropagate and step AdamW.
\end{enumerate}

\subsection{Projected DPO}
\label{sec:proj-dpo}

\warn{``Projected DPO'' is not a standard term in the literature. You must define precisely
what projection means before implementing. Three common interpretations are: (a)
\textbf{gradient projection} --- project policy gradients onto a feasible set (e.g.\
enforce a KL constraint via a Lagrange multiplier); (b) \textbf{policy projection} ---
after each gradient step, project $\pi_\theta$ back into a ball around $\pi_\text{ref}$;
(c) \textbf{subspace projection} --- restrict parameter updates to a low-rank subspace
(similar to LoRA). Without a precise definition the method cannot be reproduced or compared.}

\todo{Write a formal description of projected DPO once the variant is decided. Add the
modified loss or update rule here.}

\subsection{Hierarchical EM (Optional Extension)}
\label{sec:hier-em}

In the base Dawid--Skene model all annotators are treated as independent. A hierarchical
extension pools information across annotators of the same type by placing a shared Beta
prior on $(\alpha_j, \beta_j)$ within each group $g$:
\[
  \alpha_j \mid g(j) \sim \mathrm{Beta}(a_g, b_g), \qquad
  \beta_j  \mid g(j) \sim \mathrm{Beta}(c_g, d_g)
\]
The M-step must then update both the per-worker parameters and the group-level
hyperparameters $(a_g, b_g, c_g, d_g)$.

\warn{Hierarchical EM is significantly more complex to implement than the flat variant.
The M-step no longer has a closed form for the Beta hyperparameters; moment matching or
variational updates are typically used. Set a firm cut-off date: if not complete by that
date, drop this and present Phase~5 as the final result.}

\todo{Formalise the modified E- and M-steps, and add the expertise-level variable if used.}

\newpage

% ============================================================
\section{Implementation}

\subsection{Project Layout}

\begin{table}[H]
  \centering
  \caption{Repository modules and their roles.}
  \label{tab:layout}
  \begin{tabular}{lll}
    \toprule
    \textbf{Module} & \textbf{File} & \textbf{Role} \\
    \midrule
    Data     & \texttt{simulator.py}           & Bipartite vote simulation on HH-RLHF \\
    Data     & \texttt{prepare\_tensors.py}     & GPT-2 tokenisation $\to$ \texttt{rlhf\_tokens.pt} \\
    Data     & \texttt{dataset.py}              & Pre-tokenised PyTorch Dataset \\
    Data     & \texttt{dataset\_wo\_sft.py}     & On-the-fly tokenised Dataset \\
    Model    & \texttt{em.py}                   & \texttt{DawidSkeneEM} (E-step, M-step) \\
    Model    & \texttt{dpo.py}                  & Sequence log-probs + weighted DPO loss \\
    Model    & \texttt{sft\_model.py}           & GPT-2 policy/reference loader \\
    Model    & \texttt{dummy\_llm.py}           & Small causal LM for ablations \\
    Training & \texttt{train\_joint.py}         & Full GPT-2 joint pipeline \\
    Training & \texttt{train\_joint\_wo\_sft.py}& Ablation: dummy LM \\
    Training & \texttt{em\_only.py}             & EM baseline (no LLM priors) \\
    Utils    & \texttt{device.py}               & CUDA / MPS / CPU selection \\
    Utils    & \texttt{metrics.py}              & Golden accuracy, reward margin \\
    \bottomrule
  \end{tabular}
\end{table}

\subsection{Key Hyperparameters}

\begin{table}[H]
  \centering
  \caption{Hyperparameters for the main experimental runs.}
  \label{tab:hparams}
  \begin{tabular}{lll}
    \toprule
    \textbf{Parameter} & \textbf{Value} & \textbf{Notes} \\
    \midrule
    Tasks (prompts)       & 1,000   & First 1,000 rows of HH-RLHF train split \\
    Annotators            & 50      & 3 per task, 60 tasks per worker \\
    Batch size            & 3       & Aligned to 3 votes per prompt \\
    LR (GPT-2 pipeline)   & 1e-5    & AdamW \\
    LR (dummy LM)         & 5e-5    & \texttt{train\_joint\_wo\_sft.py} \\
    Max token length      & 128     & GPT-2 tokeniser, pad = eos \\
    EM epochs (standalone)& 15      & \texttt{em\_only.py} \\
    Training epochs (dummy)& 30     & \texttt{train\_joint\_wo\_sft.py} \\
    \bottomrule
  \end{tabular}
\end{table}

\subsection{Setup}

\begin{lstlisting}[language=bash, caption={Environment setup}]
python -m venv .venv
source .venv/bin/activate        # macOS / Linux
pip install -r requirements.txt  # torch, transformers, datasets, pandas, numpy, tqdm
\end{lstlisting}

Device selection is automatic: CUDA $>$ Apple MPS $>$ CPU (\texttt{src/utils/device.py}).

\todo{Add the specific random seed(s) and hardware spec used for the 75\% result.}

\newpage

% ============================================================
\section{Experimental Design \& Results}

The evaluation follows a staged pipeline in which each phase fixes the best configuration
found so far and adds one new component. This allows clean ablation of EM quality, DPO
variant, and SFT, alongside a robustness stress test.

\subsection{Metrics}

\paragraph{Golden dataset accuracy.}
The fraction of prompts where the predicted preference matches the simulator's hidden
\texttt{truth\_is\_A} label. Two variants are tracked:
\begin{itemize}[noitemsep]
  \item \textit{EM accuracy:} $\hat{y}_i = \mathbf{1}[\gamma_i > 0.5]$ vs.\ \texttt{truth\_is\_A}.
  \item \textit{DPO accuracy:} $\hat{y}_i = \mathbf{1}[r(x_i,y_i^A) > r(x_i,y_i^B)]$ vs.\ \texttt{truth\_is\_A}.
\end{itemize}

\paragraph{EM quality loss.}
A scalar summarising how close the inferred parameters are to ground truth:
\[
  \mathcal{L}_\text{EM} =
    \underbrace{\frac{1}{N}\sum_i \bigl(\gamma_i - \texttt{truth\_is\_A}_i\bigr)^2}_{\text{label MSE}}
  + \underbrace{\frac{1}{J}\sum_j \left(\frac{\alpha_j+\beta_j}{2} - \theta_j^\text{true}\right)^2}_{\text{worker MSE}}
\]
Used for all EM convergence plots (not a training signal --- requires simulator ground truth).

\warn{$\mathcal{L}_\text{EM}$ is an \emph{evaluation} metric only. Never use it as a loss
during training; that would leak the hidden ground truth into the model.}

\subsection{Data Setup}

The first 1,000 rows of the Anthropic HH-RLHF training split are used throughout.
\begin{itemize}[noitemsep]
  \item \textbf{Training set:} prompts used for vote simulation and policy training.
  \item \textbf{Validation set:} a held-out slice (e.g.\ 20\%) for epoch-wise monitoring.
  \item \textbf{Ground truth:} \texttt{truth\_is\_A} per prompt --- used only for metrics,
        never for training.
\end{itemize}

\warn{Fix the train/val split and random seed \emph{before} running any experiment.
Any change to the split afterwards invalidates cross-phase comparisons.}

\todo{Specify exact split sizes and the random seed used.}

% ─────────────────────────────────────────────────────────────
\subsection{Phase 1 --- Baseline: Plain DPO (No EM, No SFT, No Projection)}

\paragraph{Setup.}
Aggregate annotator votes via majority vote to obtain hard labels $\hat{y}_i^\text{maj}$.
Train a GPT-2 policy with standard DPO against a frozen pretrained GPT-2 reference
($\pi_\text{ref}$ = base GPT-2, not fine-tuned on HH-RLHF).

\paragraph{Plots.} DPO loss, DPO golden accuracy, and implicit reward margin
$\mathbb{E}[r(x,y^A)-r(x,y^B)]$, all vs.\ epoch.

\paragraph{Expected result.}
A noisy but improving baseline. Majority vote is a weak aggregation strategy; final
accuracy should plateau below all EM-assisted runs.

\warn{``No SFT'' must be defined consistently: here it means $\pi_\text{ref}$ is the
off-the-shelf pretrained GPT-2. Keep this reference fixed across all DPO phases; changing
it mid-pipeline makes comparisons invalid.}

\todo{Record the final baseline accuracy. All subsequent phases are reported relative to this.}

% ─────────────────────────────────────────────────────────────
\subsection{Phase 2 --- Offline EM + Baseline DPO}

\paragraph{Setup.}
Run Dawid--Skene EM on the vote CSV (no LLM priors, offline) to produce soft labels
$\gamma_i$. Feed these into the same DPO objective as Phase~1 (same reference, no
projection, no SFT).

\paragraph{Plots.}
\begin{itemize}[noitemsep]
  \item $\mathcal{L}_\text{EM}$ vs.\ EM iteration (one curve, scalar per step).
  \item Histogram snapshots of $\gamma_i$ at iterations 1, 5, 10, 15 --- not 1,000
        individual curves. Mean$\pm$std per worker type is also useful.
  \item $\alpha_j$ and $\beta_j$ vs.\ EM iteration, grouped by annotator type.
  \item DPO golden accuracy vs.\ epoch, overlaid with the Phase~1 curve.
\end{itemize}

\paragraph{Expected result.}
EM soft labels outperform majority vote; DPO accuracy improves vs.\ Phase~1.

\warn{Plotting 1,000 individual $\gamma_i$ trajectories is unreadable. Use summary
statistics (mean, quantiles, or histograms) to visualise the $\gamma$ distribution.}

\todo{Fill in final DPO accuracy and $\mathcal{L}_\text{EM}$ at convergence.}

% ─────────────────────────────────────────────────────────────
\subsection{Phase 3 --- Offline EM + Projected DPO + SFT}

\paragraph{Setup.}
Freeze the EM CSV from Phase~2. Replace baseline DPO with projected DPO
(see \S\ref{sec:proj-dpo}). The reference model is now a GPT-2 fine-tuned via SFT on
HH-RLHF chosen responses.

\paragraph{Plots.}
DPO loss, golden accuracy, and reward margin vs.\ epoch, overlaid with Phases~1 and~2.

\paragraph{Expected result.}
Best accuracy so far. This \textbf{fixes the DPO configuration} used in Phases~5 and~6.

\warn{Using offline (fixed) EM decouples the EM quality from DPO training --- the policy
has no feedback into the E-step. This is clean for ablation but differs from the joint
loop in \texttt{train\_joint.py}. Verify this is the intended design.}

\todo{Fill in final accuracy. Confirm this is better than Phase~2 before proceeding.}

% ─────────────────────────────────────────────────────────────
\subsection{Phase 4 --- EM Stress Testing}
\label{sec:stress}

Fix the best DPO configuration from Phase~3. Vary simulation parameters to stress EM
in isolation, and report $\mathcal{L}_\text{EM}$ at convergence. Run each factor
independently first, then combine for a worst-case regime.

\begin{table}[H]
  \centering
  \caption{Stress-test factors and their expected effect on EM quality.}
  \label{tab:stress}
  \begin{tabular}{p{3.2cm}p{3cm}p{6cm}}
    \toprule
    \textbf{Factor} & \textbf{Values} & \textbf{Expected effect} \\
    \midrule
    Scale (tasks $\uparrow$, annotators fixed)
      & 1k, 2k, 5k tasks
      & Each worker's vote count drops $\to$ weaker $\alpha_j/\beta_j$ estimates \\
    Adversarial mix (spammer + adversary fraction)
      & 30\%, 50\%, 70\%
      & Noisier votes $\to$ slower convergence, higher $\mathcal{L}_\text{EM}$ \\
    Vote sparsity (votes per task)
      & 3 (baseline), 2, 1
      & Fewer observations per prompt $\to$ less certain $\gamma_i$ \\
    Combined worst case
      & High adversary + high sparsity + fixed annotators
      & Maximum stress; used in Phase~5 \\
    \bottomrule
  \end{tabular}
\end{table}

\paragraph{Plots.}
$\mathcal{L}_\text{EM}$ at convergence vs.\ each factor (bar or line chart). Individual
$\alpha_j$/$\beta_j$ recovery curves per annotator type for a representative setting.

\warn{\textbf{Scaling direction.} Increasing tasks and annotators proportionally at fixed
sparsity gives EM \emph{more} data --- it should \emph{improve} EM, not deteriorate it.
Deterioration arises when tasks scale but annotators are \emph{held fixed} (each worker's
vote count per task falls), or when the absolute votes-per-task drops. The table above
reflects the correct interpretation (annotators fixed, tasks scale up). Resolve this
before running experiments.}

\todo{Specify exact parameter values; record $\mathcal{L}_\text{EM}$ and DPO accuracy
per setting. Identify the single worst-case configuration to carry into Phase~5.}

% ─────────────────────────────────────────────────────────────
\subsection{Phase 5 --- Robustness Evaluation}

\paragraph{Setup.}
Take the worst-case vote dataset identified in Phase~4. Evaluate two configurations:
\begin{enumerate}[noitemsep]
  \item \textbf{Best DPO only (no EM):} majority vote from the noisy dataset as hard
        labels; same DPO + SFT config from Phase~3.
  \item \textbf{EM + Best DPO:} run EM on the worst-case dataset (producing degraded
        but non-trivial $\gamma_i$); same DPO config.
\end{enumerate}

\paragraph{Claim.}
Even with degraded EM labels, configuration (2) should outperform (1), demonstrating
robustness of the full pipeline to adversarial annotator conditions.

\paragraph{Plots.}
Accuracy curves for both configurations on the same axes, with Phase~1 baseline for reference.

\warn{\textbf{Baseline sanity check.} If the worst-case dataset has, say, 70\% adversarial
workers, majority vote will be \emph{systematically inverted} --- configuration (1) may
fall \emph{below 50\%} accuracy. This makes EM look impressive but trivially so. Consider
whether a stronger no-EM baseline (e.g.\ annotator-type-weighted majority, without full EM)
would make the robustness claim more convincing.}

\todo{Fill in final accuracy for both configurations.}

% ─────────────────────────────────────────────────────────────
\subsection{Phase 6 (Optional) --- Hierarchical EM + Projected DPO + SFT}

If time permits, extend the annotator model to include an explicit expertise level
(see \S\ref{sec:hier-em}). Combined with the best DPO + SFT variant, this should yield
the highest golden accuracy across all phases.

\paragraph{Plots.}
Final accuracy from all phases on a single ``staircase'' improvement plot.

\todo{Implement hierarchical EM and fill in results. If not completed, present Phase~5 as the conclusion.}

% ─────────────────────────────────────────────────────────────
\subsection{Summary of Results}

\begin{table}[H]
  \centering
  \caption{Golden-dataset accuracy across all pipeline phases.}
  \label{tab:summary}
  \begin{tabular}{clll}
    \toprule
    \textbf{Phase} & \textbf{Configuration} & \textbf{DPO Acc.} & \textbf{$\mathcal{L}_\text{EM}$} \\
    \midrule
    1 & Plain DPO --- majority vote, no EM, no SFT            & \todo{fill} & --- \\
    2 & Offline EM + Baseline DPO                              & \todo{fill} & \todo{fill} \\
    3 & Offline EM + Projected DPO + SFT                       & \todo{fill} & same as Ph.\ 2 \\
    4 & EM stress test (worst single factor)                   & \todo{fill} & \todo{fill} \\
    4 & EM stress test (combined worst case)                   & \todo{fill} & \todo{fill} \\
    5 & Worst-case data + Best DPO, no EM                      & \todo{fill} & --- \\
    5 & Worst-case data + EM + Best DPO                        & \todo{fill} & \todo{fill} \\
    6 & Hierarchical EM + Projected DPO + SFT (optional)       & \todo{fill} & \todo{fill} \\
    \bottomrule
  \end{tabular}
\end{table}

\newpage

% ============================================================
\section{Discussion}

\subsection{Effect of LLM Priors on EM Convergence}

\todo{Discuss whether and how the policy prior changes $\gamma$ convergence relative to
the flat 50/50 prior. Does a stronger policy lead to faster / better convergence?}

\subsection{Limitations}

\begin{itemize}[noitemsep]
  \item Simulated annotators may not capture real human disagreement patterns.
  \item Only 1,000 prompts --- small scale relative to real RLHF pipelines.
  \item The E-step in \texttt{em.py} adds the global $\pi_1$ term on top of the LLM prior
        (and applies the static prior twice when no LLM prior is given), deviating from
        textbook Dawid--Skene.
\end{itemize}

\todo{Add any further limitations identified during write-up.}

\subsection{Open Design Questions}

The following concerns should be resolved before or during the demo. They are flagged here
so they do not get lost.

\begin{enumerate}

  \item \textbf{``Projected DPO'' is undefined.}
        The term does not appear in the standard literature. A precise mathematical
        definition must be written before implementation (see \S\ref{sec:proj-dpo}).
        Without it, the Phase~3 improvement cannot be attributed to projection vs.\ SFT.

  \item \textbf{Reference model consistency.}
        Phase~1 uses pretrained GPT-2 as $\pi_\text{ref}$; Phase~3 uses an SFT-tuned GPT-2.
        Changing the reference changes the implicit reward scale, making direct accuracy
        comparisons across phases potentially confounded. Consider reporting reward margins
        in addition to accuracy to make this visible.

  \item \textbf{Scaling direction in Phase~4.}
        Increasing tasks and annotators proportionally at fixed vote-sparsity should
        \emph{improve} EM, not stress it. The intended stress regime is: tasks increase,
        annotators fixed (so each worker covers more tasks with the same budget, reducing
        per-worker signal density). Confirm the simulation code reflects this.

  \item \textbf{Majority-vote baseline in Phase~5.}
        A highly adversarial dataset (e.g.\ 70\% adversary workers) will produce
        systematically inverted majority votes, driving baseline accuracy well below 50\%.
        The EM result will look strong by comparison, but the contrast is trivial.
        A more informative baseline would be a weighted combination using prior knowledge
        of annotator type fractions (without full EM).

  \item \textbf{The double-prior bug in \texttt{em.py}.}
        The E-step currently applies the global $\pi_1$ prior \emph{in addition} to the
        LLM prior, deviating from textbook Dawid--Skene. This inflates the effect of the
        global prior and may artificially advantage phases that use LLM priors. Fix before
        final experiments if possible.

  \item \textbf{$\gamma_i$ visualisation scale.}
        Plotting all 1,000 $\gamma_i$ trajectories over EM iterations is unreadable.
        Use the scalar $\mathcal{L}_\text{EM}$, histogram snapshots, or per-type
        mean $\pm$ std instead.

\end{enumerate}

\subsection{Future Work}

\begin{itemize}[noitemsep]
  \item Scale to the full HH-RLHF dataset.
  \item Replace simulated annotators with real crowdworker disagreement data.
  \item Experiment with a larger base model (e.g.\ GPT-2 medium/large or a modern open LLM).
  \item Fix the double-prior issue in \texttt{em.py} and re-evaluate.
\end{itemize}

\todo{Add any future directions specific to your course goals.}

\newpage

% ============================================================
\section{Conclusion}

This project demonstrates a principled approach to learning from noisy crowd preferences.
Dawid--Skene EM recovers per-prompt latent preferences and per-worker reliability, and
LLM priors from the current policy enrich the E-step. The resulting soft labels drive a
weighted DPO-style training objective that achieves $\approx$75\% golden-dataset accuracy
with a GPT-2 policy.

\todo{Revise once final numbers are in and key takeaways are clear.}

\newpage

% ============================================================
\bibliographystyle{plain}
\begin{thebibliography}{9}

\bibitem{dawid1979}
A.~P. Dawid and A.~M. Skene.
\newblock Maximum likelihood estimation of observer error-rates using the EM algorithm.
\newblock \textit{Applied Statistics}, 28(1):20--28, 1979.

\bibitem{rafailov2023}
R.~Rafailov, A.~Sharma, E.~Mitchell, et al.
\newblock Direct preference optimization: Your language model is secretly a reward model.
\newblock \textit{NeurIPS}, 2023.

\bibitem{bai2022}
Y.~Bai et al.
\newblock Training a helpful and harmless assistant with reinforcement learning from human feedback.
\newblock \textit{arXiv:2204.05862}, 2022.

\end{thebibliography}

\todo{Add all remaining citations: InstructGPT, original RLHF papers, crowd annotation work.}

\newpage

% ============================================================
\appendix
\section{Running the Code}

All data scripts must be run from \texttt{src/data/} and training scripts from
\texttt{src/training/} so relative \texttt{../../data/} paths resolve correctly.

\begin{table}[H]
  \centering
  \caption{Run commands and their outputs.}
  \begin{tabular}{lll}
    \toprule
    \textbf{Step} & \textbf{Command} & \textbf{Output} \\
    \midrule
    1. Simulate votes  & \texttt{cd src/data \&\& python simulator.py}            & \texttt{data/processed/*.csv} \\
    2. Pre-tokenise    & \texttt{cd src/data \&\& python prepare\_tensors.py}      & \texttt{data/tokenized/*.pt} \\
    3. Full training   & \texttt{cd src/training \&\& python train\_joint.py}      & Metrics to terminal \\
    3. Ablation        & \texttt{cd src/training \&\& python train\_joint\_wo\_sft.py} & Metrics to terminal \\
    EM only            & \texttt{cd src/training \&\& python em\_only.py}          & Accuracy + worker MAE \\
    Loss smoke test    & \texttt{cd src/training \&\& python test\_dpo.py}         & Forward/backward check \\
    \bottomrule
  \end{tabular}
\end{table}

\section{Known Implementation Note}

In \texttt{src/models/em.py}, the E-step adds both the LLM prior term and the global
$\pi_1$ term to each branch. When \texttt{llm\_priors} is \texttt{None}, the static 50/50
prior is effectively applied twice in log-space. When LLM priors are provided, the global
$\pi_1$ term is still added on top. This deviates from the textbook Dawid--Skene
formulation and should be reviewed before publication.

\end{document}
